\subsection{Macaulay rings}

DEFINITION 3.--- A ring A is said to be Macaulay, or to be a Macaulay ring, if it is Noetherian and the A-module A is Macaulay.

Examples. - 1) Every Artinian ring is a Macaulay ring (no. 1, example 1). 2) A Macaulay ring has no embedded associated prime ideals (no. 2, prop. 2). Conversely, let A be a Noetherian ring of dimension ≤ 1 with no embedded associated prime ideals; for every nonempty finite strongly secant subset S of A, the A-module A/SA has dimension ≤ 0, hence is Macaulay (no. 1, example 1); consequently A is a Macaulay ring (no. 4, th. 2). In particular, a reduced Noetherian ring of dimension ≤ 1 is a Macaulay ring.

\begin{enumerate}
\def\labelenumi{\arabic{enumi})}
\setcounter{enumi}{2}
\item
  A normal Noetherian ring of dimension \(\leqslant 2\) is a Macaulay ring (\(\S 1, n^{\circ} 10\), text preceding Th. 4). Conversely, let \(A\) be a Macaulay ring whose local ring \(A_{\mathfrak{p}}\) at every prime ideal \(\mathfrak{p}\) of height \(\leqslant 1\) is integrally closed; then A is normal (\(\S 1, n^{\circ} 10\), cf. th. 4).
\item
  If A is a Macaulay ring, then so is \(S^{-1}A\) for every multiplicative subset S of A (\(n^{\circ}\) 1, remark). Conversely, if the ring \(A_{m}\) is a Macaulay ring for every maximal ideal m of A, then the ring A is Macaulay (\(n^{\circ}\) 1, def. 1).
\item
  Let A be a Noetherian ring and J an ideal of \(A\). For A/J to be a Macaulay ring, it is necessary and sufficient that it be an \(A\)-Macaulay module (no. 1, example 4).
\item
  Let A be a Noetherian local ring and J an ideal of A generated by an A-regular sequence. For A/J to be a Macaulay ring, it is necessary and sufficient that A be one (example 5 and prop. 4 of no. 3).
\item
  For a Noetherian local ring A to be Macaulay, it is necessary and sufficient that it possess an ideal of definition generated by an A-regular sequence: this follows from prop. 3 of no. 3, and from the fact that an A-regular sequence of elements of m\_A generates an ideal of definition if and only if it has length dim(A) (VIII, § 3, no. 2, th. 1 and cor. of prop. 3). In particular, every regular Noetherian local ring is a Macaulay ring (VIII, § 5, no. 2, th. 1). More generally, the quotient of a regular Noetherian local ring A by an ideal generated by an A-regular sequence is a Macaulay ring (example 6).
\end{enumerate}

PROPOSITION 7. -- Let A be a Noetherian ring. The following conditions are equivalent:

\begin{enumerate}
\def\labelenumi{(\roman{enumi})}
\item
  A is a Macaulay ring;
\item
  for every closed subset F of Spec(A), one has \(\mathrm{prof}_{\mathrm{F}}(\mathrm{A}) = \mathrm{codim}(\mathrm{F})\) ;
\item
  every ideal J of A contains an A-regular sequence of length ht(J) ;
\end{enumerate}

(iii') every maximal ideal m of A contains an A-regular sequence of length ht(m) ;

\begin{enumerate}
\def\labelenumi{(\roman{enumi})}
\setcounter{enumi}{3}
\tightlist
\item
  for every ideal J of A, one has \(\mathrm{Ext}_A^i (\mathrm{A} / \mathrm{J},\mathrm{A}) = 0\) for \(i <   \mathrm{ht}(\mathrm{J})\)
\end{enumerate}

(iv') for every maximal ideal m of A, one has Ext \(_{A}^{i}\) (A/m, A) = 0 for i \textless{} ht(m);

\begin{enumerate}
\def\labelenumi{(\alph{enumi})}
\setcounter{enumi}{21}
\tightlist
\item
  for every prime ideal p of A and every ideal J of \(A_{p}\) generated by a maximal secant sequence for \(A_{p}\) , one has \(e_{\mathrm{J}}(\mathrm{A}_{\mathfrak{p}})=\mathrm{long}(\mathrm{A}_{\mathfrak{p}}/\mathrm{JA}_{\mathfrak{p}})\) ;
\end{enumerate}

(v') for every maximal ideal m of A, there exists an ideal J of \(A_{m}\) , generated by a maximal secant sequence for \(A_{m}\) , satisfying \(e_{\mathrm{J}}(A_{\mathrm{m}})=\mathrm{long}(\mathrm{A}_{\mathrm{m}}/\mathrm{JA}_{\mathrm{m}})\) ;

\begin{enumerate}
\def\labelenumi{(\roman{enumi})}
\setcounter{enumi}{5}
\tightlist
\item
  (Macaulay-Cohen criterion) for every finite subset S of A such that the ideal G generated by S has height Card(S), the A-module A/G has no embedded associated prime ideals.
\end{enumerate}

The equivalence of (i) and (ii) follows from no. 1, cor. of prop. 1. By th. 2 of § 1, no. 4, and the definition of depth, conditions (iii) and (iv) (resp. (iii') and (iv')) mean that one has \(\mathrm{prof}_{\mathrm{A}}(\mathrm{J};\mathrm{A})\geqslant \mathrm{ht}(\mathrm{J})\) for every ideal (resp. every maximal ideal) J of A. One therefore has

\[
(i) \Leftrightarrow (ii) \Rightarrow (iii) \Leftrightarrow (iv) \Rightarrow (iii') \Leftrightarrow (iv').
\]

But (iv') implies, for every maximal ideal m of A, \(\operatorname{Ext}_{A_{m}}^{i}(\kappa(m), A_{m}) = 0\) for \(i < \dim(A_{m})\) , whence \(\operatorname{prof}(A_{m}) \geqslant \dim(A_{m})\) , so that A is a Macaulay ring.

The equivalence of (i), (v) and (v') follows from th. 1 of no. 3, and that of (i) and (vi) from th. 2 of no. 4.

